% Q2: Motional EMF — Sliding bar on parallel rails with resistor
\documentclass[border=10pt]{standalone}
\usepackage[american voltages, american currents]{circuitikz}
\usepackage{amsmath}
\renewcommand{\familydefault}{\sfdefault}

\begin{document}
\begin{circuitikz}[
    line width=0.8pt,
    every node/.style={font=\sffamily}
]

% Top rail
\draw[thick] (0,4) -- (7,4);
% Bottom rail
\draw[thick] (0,0) -- (7,0);

% Resistor on left side closing the circuit
\draw (0,4) to[R, l=$R$] (0,0);

% Sliding bar (vertical, at x=5)
\draw[very thick, blue!70!black] (5,0) -- (5,4);

% Bar length label
\draw[<->, >=stealth] (5.4,0) -- (5.4,4)
    node[midway, right] {$\ell$};

% Velocity arrow
\draw[->, >=stealth, very thick, red!70!black] (5,2) -- (6.5,2)
    node[above] {$\mathbf{v}$};

% B field symbols (into page — X marks)
\foreach \x in {1.5, 2.5, 3.5} {
    \foreach \y in {1.0, 2.0, 3.0} {
        \node[gray!60!black] at (\x, \y) {$\times$};
    }
}

% B field label
\node[gray!60!black] at (2.5, -0.6) {$\mathbf{B}$ (into page)};

% Current direction arrow (counterclockwise in this view)
\draw[->, >=stealth, thick, green!50!black] (1.0, 3.6) -- (3.5, 3.6)
    node[midway, above] {$I$};

% Rail terminations (dots at sliding contact)
\fill (5,4) circle (2pt);
\fill (5,0) circle (2pt);

\end{circuitikz}
\end{document}
